% This document is compiled using pdfLaTeX
% You can switch XeLaTeX/pdfLaTeX/LaTeX/LuaLaTeX in Settings

\documentclass{article}
\usepackage{ctex}
\usepackage{amsmath} % 确保引入了此宏包以支持 aligned

\title{因子日历2026}

\date{\today}

\begin{document}

\maketitle

\section{加权偏度(高频因子-收益分布类)}

\subsection{因子定义}
\begin{equation}
    \text{加权偏度} = \frac{ \sum_{t=1}^T w_t (\mathrm{close}_t - \overline{\mathrm{close}})^3}{\mathrm{close}_{\sigma}^3}
\end{equation}

\subsection{变量定义}
\begin{itemize}
    \item ① \( w_t = \frac{\mathrm{vol}_t}{VOL} \) 为正比于成交量的权重；
    \item ② \( \overline{\mathrm{close}} \) 为个股收盘价均值；
    \item ③ \( \mathrm{close}_{\sigma} \) 为个股收盘价标准差；
    \item ④ \( T \) 为划分时间段的个数，如日内 5 分钟频率的 \( T=48 \)。
\end{itemize}

\subsection{说明}
上述加权偏度是以成交量加权计算的价格分布的偏离程度，在计算偏度时，给成交量较高的时间段更大的权重；与未来收益负相关，负偏态越明显，个股成交集中在价格相对高位的水平，容易造成股价的高估。

\subsection{参考文献}
\begin{enumerate}
    \item 覃川桃, 郑起. 2020. 高频因子（八）：高位成交因子——从量价匹配说起. 长江证券.
\end{enumerate}

\section{“朝没晨雾”因子(高频因子-量价相关性类)}

\subsection{因子定义}

1. 对股票 \( i \)，用每日 1 分钟收盘价计算得到第 \( t \) 分钟收益率序列 \( \mathrm{ret}_t \)；用每日 1 分钟成交量计算得到 1 分钟增量成交量 \( \mathrm{voldiff}_t \)：
   \begin{equation}
       \mathrm{ret}_t = \frac{\mathrm{close}_t}{\mathrm{close}_{t-1}} - 1
   \end{equation}
   \begin{equation}
       \mathrm{voldiff}_t = \mathrm{volume}_t - \mathrm{volume}_{t-1}
   \end{equation}

2. 对每天第 6 分钟至第 240 分钟的上述数据进行带截距项的最小二乘回归：
   \begin{equation}
       \mathrm{ret}_t = \alpha_0 + \alpha_1 \mathrm{voldiff}_t + \alpha_2 \mathrm{voldiff}_{t-1} + \cdots + \alpha_6 \mathrm{voldiff}_{t-5} + \varepsilon_t
   \end{equation}

3. 记上述回归得到的第 \( t-1 \)、\( t-2 \)、\( t-3 \)、\( t-4 \)、\( t-5 \) 分钟的增量成交量的回归系数的 \( t \) 值分别为 \( t_1 \)、\( t_2 \)、\( t_3 \)、\( t_4 \)、\( t_5 \)，计算这 5 个 \( t \) 值的标准差，即为股票 \( i \) 在 \( T \) 日的“日朝没晨雾”因子；

4. 每月末，计算过去 20 个交易日的“日朝没晨雾”因子均值，即得到“朝没晨雾”因子。

\subsection{说明}
“朝没晨雾”因子衡量了每分钟的前 5 分钟的信息到来的平稳程度，与未来收益负相关；因子取值越小，表示该股票短时间内流入的信息越平稳，越没有突然到来的信息，表明该股票相对冷门，投资者对其的关注较少并且相对理智，因此其当前价格可能相对被低估，未来可能会产生较高收益。

\subsection{参考文献}
\begin{enumerate}
    \item 曹春晓. 推动个股价格变化的因素分解与“花隆林间”因子, 2023, 方正证券.
\end{enumerate}

\section{早盘主动净流入率稳定性(量价因子改进)}

\subsection{因子定义}
剔除当月早盘振幅最高的 \(1/5\) 标的，对剩余标的计算如下因子：
\begin{equation}
    \text{早盘主动净流入率稳定性因子}=\frac{\mathrm{mean}(\text{当月每日早盘主动净流入率}_{i,t})}{\mathrm{std}(\text{当月每日早盘主动净流入率}_{i,t})}
\end{equation}
其中，
\begin{equation}
    \text{早盘主动净流入率}_{i,t} = \frac{\text{开盘10:00前净主动买入额}_{i,t}}{\text{10:00前总成交额}_{i,t}}
\end{equation}

\subsection{说明}
该因子衡量了早盘主动资金流入的稳定性，可用于筛选主动资金近期具有稳定偏好的标的；使用早盘振幅是为了剔除非稳定性资金干扰较高的标的。

\subsection{参考文献}
\begin{enumerate}
    \item 王琦. 2023. JASON's alpha: 基本面+量价复合策略. 东北证券.
\end{enumerate}

\section{成交量支撑区域下限与收盘价差异（高频因子-成交分布类）}

\subsection{因子定义}

1. 计算同价成交量：将日内相同分钟收盘价的成交量累加至一起，得到在当前价格上的成交量总和；

2. 将成交量累计最大的价格定义为该股当日的成交量支撑点 \( VolumeSupportPrice \)，包含该支撑点的附近区域定义为成交量支撑区域 \( VolumeSupportArea \)；

3. 逐步计算 \( \mathrm{VSP} \) 周围的成交量累计值（价格由近及远，成交量由大及小），该累计值与全天成交量总和的比值超 50\% 的最小区域，即为成交量支撑区域 \( VSA \)，该区域上限价格为 \( \mathrm{VSA\_High} \)，下限价格为 \( \mathrm{VSA\_Low} \)；

4. 计算下限价格 \( \mathrm{VSA\_Low} \) 与当日收盘价之间的差异即得到 \( \mathrm{vsa\_ratio} \)：

\subsection{说明}
成交量支撑区域是判断个股日内公允价值区域的重要依据，过高/过低于该区域的价格将难以维持；若下限价格与当日收盘价之间差异较大，说明该股收盘后处于异常推动状态，当日大幅跌穿公允价值区域，其未来预计会回到公允价值区域，出现回弹上涨。

\subsection{参考文献}
\begin{enumerate}
    \item 郑兆磊. 2022. 成交量分布中的 Alpha. 兴业证券.
\end{enumerate}

\section{大的上行跳跃波动率(高频因子-波动跳跃类)}

\subsection{因子定义}
\begin{equation}
    \mathrm{RVLJP}_{\gamma,t} = \min\left(\mathrm{RVJP}_t, \sum_{i=1}^n r_{i,t}^2 \mathrm{I}_{\{r_{i,t} > \gamma\}}\right)
\end{equation}
\begin{equation}
    \mathrm{RVJP}_t = \max\left(\mathrm{RS}_t^+ - \frac{1}{2} \hat{IV}_t, 0\right)
\end{equation}
\begin{equation}
    \gamma = \alpha \Delta_n^{0.49} \sqrt{\hat{IV}_t}
\end{equation}

\subsection{变量定义}
\begin{itemize}
    \item ① \( \mathrm{RVJP}_t \) 为交易日 \( t \) 的上行跳跃波动率，计算细节可参考相关日历页；
    \item ② \( \gamma \) 为区分大、小跳跃的阈值，\( \Delta_n \) 为日内分钟数据个数的倒数，如 5 分钟频率的 \( \Delta_n = 1/48 \)，经验参数 \( \alpha = 4 \)；
    \item ③ \( r_{i,t} \) 为某股票交易日 \( t \) 内的分钟对数收益率序列，如 5 分钟对数收益率序列。
\end{itemize}

\subsection{说明}
跳跃波动率可以根据阈值进一步分解为大、小两个部分，即用于区分信息冲击导致的大跳跃波动和无法用信息冲击解释的小跳跃波动。上行波动率也可以做上述划分；由于大信息冲击更有可能改变投资者对股票基本面的预期，所以大跳跃波动通常比小跳跃波动有更强的收益预测能力。

\subsection{参考文献}
\begin{enumerate}
    \item Duong, D., \& Swanson, N. R. (2015). \textit{Empirical evidence on the importance of aggregation, asymmetry, and jumps for volatility prediction}. Journal of Econometrics, 187, 606-621.
    \item Bo Yu, Bruce Mizrach, Norman R. Swanson. (2020). \textit{New Evidence of the Marginal Predictive Content of Small and Large Jumps in the Cross-Section}. Econometrics, MDPI, 8(2), 1-52.
\end{enumerate}

\section{亏损卖出倾向(行为金融因子-处置效应)}

\subsection{因子定义}

1. 计算 \( t \) 时刻收盘价 \( \mathrm{Close}_t \) 与前 \( n \) 天 VWAP 均价 \( P_{t-n} \) 之间的涨跌幅，取涨跌幅为负的部分作为亏损收益率：
\begin{equation}
    \mathrm{loss}_t = \frac{\mathrm{Close}_t - P_{t-n}}{\mathrm{Close}_t} \cdot \mathbf{1}_{\{\mathrm{Close}_t < P_{t-n}\}}
\end{equation}

2. 以换手率 \( V_t \) 为权重加权亏损收益率得到资本亏损量即为亏损卖出倾向因子 \( \mathrm{Loss} \)：
\begin{equation}
    \mathrm{Loss}_t = \sum_{n=0}^{N-1} w_{t-n} \cdot \mathrm{loss}_{t-n}
\end{equation}
\begin{equation}
    w_{t-n} = \frac{1}{k} V_{t-n} \prod_{s=1}^{n-1} (1 - V_{t-n+s})
\end{equation}
\begin{equation}
    k = \sum_{n=1}^{T} V_{t-n} \prod_{s=1}^{n-1} (1 - V_{t-n+s})
\end{equation}

\subsection{说明}
亏损卖出倾向因子衡量了大部分投资者“未实现”的亏损，因子值越小，投资者未实现亏损越大，卖出意愿越强，股票出现超卖，股价偏低，未来股价更有可能上升，与未来收益负相关。

\subsection{参考文献}
\begin{enumerate}
    \item 陈升锐, 姚紫薇. 2025. 处置效应与 V 型处置效应在量化选股中的应用. 中信建投.
    \item An, L. (2016). \textit{Asset Pricing When Traders Sell Extreme Winners and Losers}. The Review of Financial Studies, 29(3), 823-861.
\end{enumerate}

\section{散户羊群效应(高毅因子-资金流类)}

\subsection{因子定义}
\begin{equation}
    \mathrm{RankCorr}(R_t, S_{t+1})
\end{equation}

\subsection{变量定义}
\begin{itemize}
    \item ① \(R_t\) 为过去 20 个交易日的日度收益率序列，用于表示市场涨跌；
    \item ② \(S_{t+1}\) 为小单（<4万元）过去 20 个交易日周期的资金净流入序列（前置 1 期），用于表示散户行为；
    \item ③ 使用 \(\frac{\mathrm{close}_t}{\mathrm{open}_t} - 1\) 日内收益，使用小单非主动净流入（开盘至 10 点之间），可以提升原始散户羊群效应的表现。
\end{itemize}

\subsection{说明}
衡量了散户追涨杀跌的程度，与股票未来收益负相关，即股票的散户羊群效应越高，未来预期收益越差。

\subsection{参考文献}
\begin{enumerate}
    \item 魏建榕, 盛少成. 2022. 新型因子：资金流动力学与散户羊群效应. 开源证券.
    \item 魏建榕, 盛少成. 2022. 大小单资金流 alpha 研究 2.0: 变量精筛与高频测算. 开源证券.
\end{enumerate}

\section{买单非流动性(高频因子-流动性类)}

\subsection{因子定义}
\begin{equation}
    r_{i,t} = \alpha + \beta_1 * S_{i,t} + \beta_2 * B_{i,t} + \epsilon_{it}
\end{equation}

\subsection{变量定义}
\begin{itemize}
    \item ① \( \beta_1 \) 为卖出非流动性系数；
    \item ② \( \beta_2 \) 为买入非流动性系数；
    \item ③ \( S_{i,t} \) 为股票 \( i \) 在 \( t \) 时间区间内的主动卖出金额；
    \item ④ \( B_{i,t} \) 为股票 \( i \) 在 \( t \) 时间区间内的主动性买入金额。
\end{itemize}

\subsection{说明}
买单非流动性衡量了高频数据下主动买入的交易金额对于股票价格变动的影响；买单非流动性对收益率的预测效果不及卖单非流动性，主要是由于投资者存在亏损厌恶的心理。

\subsection{参考文献}
\begin{enumerate}
    \item Brennan, M. J., Chordia, T., Subrahmanyam, A., \& Tong, Q. (2012). \textit{Sell-order liquidity and the cross-section of expected stock returns}. Journal of Financial Economics, 105(3), 523–541.
    \item 朱剑涛. 2017. 技术类新 Alpha 因子的批量测试. 东方证券.
\end{enumerate}

\section{年度搜索比 (downstream)(图谱网络-动量溢出)}

\subsection{因子定义}
\begin{equation}
    f_{ij}^t = \frac{\sum_{d=1}^{365} (\text{Daily Searches for } j \text{ after } i)_d}{\sum_{d=1}^{365} (\text{Daily Searches for } i \text{ and then any firm } j \neq i)_d}
\end{equation}
\begin{equation}
    R_{i,t} = \frac{\sum_{j=1}^N f_{ij}^t \times \mathrm{Ret}_j}{\sum_{j=1}^N f_{ij}^t}
\end{equation}

\subsection{变量定义}
\begin{itemize}
    \item ① \( f_{ij}^t \) 为年度搜索比、\( \mathrm{Ret}_j \) 为月度涨跌幅；
    \item ② 原始数据为从美国证券交易委员会（SEC）的 EDGAR 网站上获取的用户对公司信息的搜索流量；
    \item ③ 只保留同一天内至少搜索了两个不同公司（基于 CIK 编号识别的不同公司）的用户记录；
    \item ④ 剔除同一天内，用户对同一公司的连续搜索，对每个公司的搜索只计数一次；
    \item ⑤ 限制每个用户每天下载的独特公司数量不超过 50 家，以减少批量下载者的影响；
    \item ⑥ 仅考虑特定类型文件的页面浏览量，如 10-K、10-Q、8-K、14-A、S-1 等报告。
\end{itemize}

\subsection{说明}
年度搜索比例反映了投资者在搜索引擎中 \( i \) 之后，进一步搜索公司 \( j \) 的频率，借助投资者搜索行为捕捉投资者认为这两家公司之间存在经济关联的关联程度；投资者搜索行为有助于识别可能未被传统方法发现但具有密切经济联系的企业，而这些联系相比行业分类等传统关联对于解释股票回报、估值倍数等财务指标的变化更加有效。

\subsection{参考文献}
\begin{enumerate}
    \item Lee, Charles M. C., Ma, Paul, \& Wang, Charles C. Y. (2015). \textit{Search-Based Peer Firms: Aggregating Investor Perceptions Through Internet Co-Searches}. Journal of Financial Economics, 116(2), 410-431.
\end{enumerate}

\section{修正的日内反转(高频因子-动量反转类)}

\subsection{因子定义}

1. 对股票 \( i \)，计算 \( t \) 日的日内收益率 = \( t \) 日收盘价 / \( t \) 日开盘价 - 1；

2. 每月末计算当月的日内收益率和日内收益率的波动率，即月末最近 20 天的日内收益率序列的均值和标准差；

3. 对股票 \( i \)，若上述日内收益率的波动率小于横截面均值，则将其当月日内收益率取相反数（未来更可能发生动量效应），进而实现了反转的修正。

\subsection{变量定义}
\begin{itemize}
    \item ① 上述因子是借助波动率对反转做修正，也可借助换手率来做修正，即换手率变化量（\( t \) 日换手率与 \( t-1 \) 日换手率的差值）低于横截面均值，未来更可能发生动量；
    \item ② 可将日内收益率替换成日间收益率，对日间反转做修正。
\end{itemize}

\subsection{说明}
提高投资者对股价未来趋势判断的“可知性”，有助于更清晰的判断股票未来是否发生反转效应还是动量效应。波动率和换手率指标可代表股票的“可知性”，低波动（股价走势稳定）或低换手（投资分歧降低）的股票可知性更高，未来更可能发生动量效应；反之，可知性更低，未来更可能发生反转。通过将预判的动量做“翻转”，即可提高反转因子的表现。

\subsection{参考文献}
\begin{enumerate}
    \item 曹春晓. 2022. 个股动量效应的识别及“球状硬币”因子构建. 方正证券.
    \item Moskowitz, T. J. (2021). \textit{Asset pricing and sports betting}. Journal of Finance, Forthcoming.
\end{enumerate}
\end{document}